\clearpage
\section{대칭다항식을 실제로 바꾸는 방법}
\label{sec:symmetric-reduction-algorithm}

\begin{learninggoal}
사전식 최고항 제거 알고리즘으로 대칭다항식을 기본대칭식에 관한 다항식으로 변환한다. 변수 수와 차수를 낮추는 재귀적 방법을 이해하고, 대칭 유리함수의 고정체를 다시 확인한다.
\end{learninggoal}

기본정리의 증명은 그대로 계산 알고리즘이 된다. 대칭다항식 $f$의 사전식 최고항이
\[
  cT_1^{\lambda_1}\cdots T_n^{\lambda_n},
  \qquad
  \lambda_1\ge\cdots\ge\lambda_n
\]
이면
\[
  c\,s_1^{\lambda_1-\lambda_2}
  s_2^{\lambda_2-\lambda_3}\cdots s_n^{\lambda_n}
\]
을 빼서 최고항을 없앤다. 나머지에 같은 절차를 반복하면 유한 단계 뒤에 끝난다.

\begin{example}[제곱합]
두 변수에서
\[
  f=T_1^2+T_2^2
\]
의 최고항은 $T_1^2$이고 지수벡터는 $(2,0)$이다. 대응하는 기본대칭식 단항식은 $s_1^2$이다. 그런데
\[
  s_1^2=(T_1+T_2)^2
  =T_1^2+2T_1T_2+T_2^2
\]
이므로
\[
  T_1^2+T_2^2=s_1^2-2s_2.
\]
이 항등식은 $2$가 단위원이라는 가정 없이 임의의 가환환에서 성립한다.
\end{example}

\begin{example}[세 변수의 세제곱합]
\label{ex:cubic-power-sum}
\[
  p_3=T_1^3+T_2^3+T_3^3
\]
의 최고항은 $T_1^3$이다. 먼저 $s_1^3$을 빼고, 생기는 혼합항을 차례로 정리하면
\[
  p_3=s_1^3-3s_1s_2+3s_3.
\]
직접 전개하여 확인할 수도 있다. 이 식을 다항식
\[
  X^3-s_1X^2+s_2X-s_3
\]
의 세 근에 적용하면 근의 세제곱합이 계수만으로 결정됨을 알 수 있다.
\end{example}

\begin{example}[쌍곱의 제곱합]
세 변수에서
\[
  q=T_1^2T_2^2+T_1^2T_3^2+T_2^2T_3^2
\]
를 생각하자. $s_2^2$을 전개하면
\[
  s_2^2=q+2T_1T_2T_3(T_1+T_2+T_3)
  =q+2s_1s_3.
\]
따라서
\[
  q=s_2^2-2s_1s_3.
\]
\end{example}

\subsection{변수 수와 차수를 낮추는 재귀법}

기본정리의 다른 증명 방식은 다음 재귀를 제안한다. $f(T_1,\dots,T_n)$이 대칭다항식이라 하자.

먼저 $T_n=0$을 대입한다. 그러면
\[
  f(T_1,\dots,T_{n-1},0)
\]
는 $n-1$개의 변수에 관한 대칭다항식이다. 귀납적으로 이를
\[
  f_0(s'_1,\dots,s'_{n-1})
\]
로 표현한다. 여기서 $s'_j$는 $T_1,\dots,T_{n-1}$의 기본대칭식이다.

이제 같은 다항식 $f_0$에 $n$변수 기본대칭식 $s_1,\dots,s_{n-1}$을 대입한다. 차이
\[
  f-f_0(s_1,\dots,s_{n-1})
\]
는 $T_n=0$일 때 사라지므로 $T_n$으로 나누어진다. 대칭성 때문에 모든 $T_i$로 나누어지고, 따라서
\[
  s_n=T_1\cdots T_n
\]
으로 나누어진다. 즉
\[
  f=f_0(s_1,\dots,s_{n-1})+s_ng
\]
인 더 낮은 차수의 대칭다항식 $g$가 존재한다. $g$에 같은 과정을 반복한다.

\begin{example}
세 변수의 대칭다항식
\[
  f=\sum_{i\ne j}T_i^2T_j
\]
를 생각하자. $T_3=0$을 대입하면
\[
  T_1^2T_2+T_2^2T_1
  =T_1T_2(T_1+T_2)
  =s'_1s'_2.
\]
따라서 첫 후보는 $s_1s_2$이다. 전개하면
\[
  s_1s_2
  =\sum_{i\ne j}T_i^2T_j+3T_1T_2T_3
  =f+3s_3.
\]
그러므로
\[
  f=s_1s_2-3s_3.
\]
\end{example}

\subsection{가중차수}

기본대칭식 $s_j$에 가중치 $j$를 준다. 즉 단항식
\[
  s_1^{a_1}\cdots s_n^{a_n}
\]
의 가중차수를
\[
  a_1+2a_2+\cdots+na_n
\]
로 정의한다. 각 $s_j$는 원래 변수들에 관해 $j$차 동차식이므로, $m$차 동차 대칭다항식은 기본대칭식으로 표현했을 때도 가중차수 $m$인 항들만 갖는다. 이 관찰은 계산 중 가능한 항을 크게 줄여 준다.

\begin{example}
세 변수에서 $4$차 동차 대칭다항식은
\[
  s_1^4,\quad s_1^2s_2,\quad s_2^2,\quad s_1s_3
\]
의 $R$-선형결합으로 표현된다. $s_4$는 세 변수에서는 존재하지 않는다. 따라서 계수 네 개만 정하면 된다.
\end{example}

\subsection{대칭 유리함수}

\begin{theorem}[대칭 유리함수의 기본정리]
$k$가 체이면
\[
  k(T_1,\dots,T_n)^{S_n}
  =k(s_1,\dots,s_n).
\]
즉 모든 대칭 유리함수는 기본대칭식들의 유리함수로 유일하게 표현된다.
\end{theorem}

\begin{proof}
포함
\[
  k(s_1,\dots,s_n)\subseteq k(T_1,\dots,T_n)^{S_n}
\]
은 자명하다. 반대로 대칭 유리함수 $q=f/g$를 택한다. 분모를 모든 순열로 옮긴 것들의 곱을 이용하여
\[
  D=\prod_{\pi\in S_n}\pi(g)
\]
라 두면 $D$는 대칭다항식이다. 또한
\[
  qD=f\prod_{\pi\ne\id}\pi(g)
\]
도 대칭다항식이다. 기본정리에 의해 분자와 분모 모두 $s_1,\dots,s_n$의 다항식이므로 $q$는 이들의 유리함수이다. 유일성은 $s_1,\dots,s_n$의 대수적 독립성에서 따른다.
\end{proof}

\begin{corollary}
\[
  [k(T_1,\dots,T_n):k(s_1,\dots,s_n)]=n!
\]
이고 이 확장은 $S_n$을 갈루아군으로 갖는다.
\end{corollary}

\begin{proof}
앞 절의 자유기저를 분수체로 올리면 차수가 $n!$ 이하이다. 한편 변수 순열이 서로 다른 $n!$개의 $k(s_1,\dots,s_n)$-자동동형을 주므로 자동동형 수에 의한 갈루아 판정에서 차수는 정확히 $n!$이다.
\end{proof}

\begin{keyidea}
대칭식은 근들의 이름에 의존하지 않는 양이다. 기본대칭식 정리는 그런 모든 양을 다항식의 계수로 다시 쓸 수 있다고 말한다. 판별식이 근의 차이로 정의되면서도 계수만으로 계산되는 이유가 바로 여기에 있다.
\end{keyidea}

\begin{checkpoint}
\begin{enumerate}[label=(\alph*)]
  \item $T_1^2+T_2^2+T_3^2$를 $s_1,s_2,s_3$로 나타내라.
  \item $\sum_{i\ne j}T_i^2T_j=s_1s_2-3s_3$을 직접 전개하여 확인하라.
  \item 대칭 유리함수의 증명에서 분모 $D$가 왜 대칭인지 설명하라.
\end{enumerate}
\end{checkpoint}
